\section{Theoretical Foundations}

\subsection{Control-system attachment is prior theory, not ARC novelty}

Bowlby's attachment system was formulated as goal directed and feedback regulated \parencite{bowlby1982attachment}. Kobak and colleagues later used control theory to analyze attachment strategies and emotion regulation during mother--teen problem solving \parencite{kobak1993control}. Petters reviewed the control-system lineage and its computational implications \parencite{petters2019control}. ARC adopts this lineage rather than claiming to originate it.

\subsection{Computational and dynamical precedents}

Stevens and Zhang modeled infant attachment as a dynamic system with explicit state evolution and recovery-like behavior \parencite{stevens2009dynamic}. Gagliardi implemented attachment as a multidimensional control system with representational set-goals and agent interactions \parencite{gagliardi2022control}. Olar and colleagues used an adaptive network model for anxious--avoidant relationship dynamics \parencite{olar2026relationship}. These precedents show that targets, state updating, feedback, and dyadic dynamics are already computationally plausible.

Their existence changes the novelty question. ARC should not be evaluated by whether it can reproduce another simulated attachment trajectory. It should be evaluated by whether its proposed variables can be measured repeatedly in people and whether the resulting temporal model adds information beyond static attachment scores and simpler time-series baselines.

\subsection{Risk regulation, responsiveness, and coregulation}

Adult relationship regulation also has established psychological precedents. Risk-regulation theory describes how dependence and perceived acceptance shape decisions to seek connection or protect the self \parencite{murray2006risk}. Coregulation and partner-buffering accounts emphasize that adult attachment, distress, separation, and recovery unfold through coupled interpersonal regulation rather than isolated intrapersonal control \parencite{sbarra2008coregulation,simpson2014buffering,overall2022buffering}. ARC therefore treats partner feedback as a measured component of the process rather than background noise.

\subsection{Closeness and autonomy are not assumed to be opposites}

ARC distinguishes experienced autonomy from interpersonal distance. Close-relationship research shows that autonomy and relatedness can be jointly satisfied rather than occupying opposite ends of one continuum \parencite{kluwer2020autonomy}. This supports a possible multidimensional relational-state representation while stopping short of establishing which coordinates ARC requires. The scalar closeness pilot therefore remains the first test; a closeness--autonomy extension must earn its added complexity empirically.

\subsection{Positive security as a candidate core construct}

Secure-base representations encode expectations that support can be available and effective and that exploration or constructive activity can resume \parencite{waters2006workingmodels,waters2019securebase,waters2025securebase}. This provides a strong theoretical reason to test positive secure-base representation, \Sec, as distinct from simply reversing \Anx\ and \Avo. It does not, however, establish that \Sec\ is necessary for the ARC architecture or that it predicts recovery beyond dimensional insecurity. ARC therefore prespecifies \Sec\ as a candidate core construct rather than a required component. Recovery is treated as an observable dynamic outcome against which this candidate can earn or fail to earn incremental validity.

\subsection{Candidate moderators rather than additional core constructs}

Mentalized affectivity concerns identifying, processing, and expressing affective states \parencite{greenberg2017mentalized}. Emotion-regulation flexibility concerns context-sensitive adaptation \parencite{mosannenzadeh2024flexibility}. Both are relevant to ARC, but their role is deliberately limited. Mentalized affectivity is a candidate moderator of relational-state inference quality; flexibility is a candidate moderator of context- and feedback-sensitive policy adaptation. Neither is required for the core feedback loop to be defined.

\subsection{The empirical gap}

The remaining gap is a measurement problem: can the sequence from perceived and desired relational conditions to discrepancy, action, partner feedback, and recovery be recovered in adult romantic relationships using intensive longitudinal data? ARC is organized around this question.
